\chapter{Contexto y Justificación}
\label{cap:02_contexto_justificacion}

El desarrollo de videojuegos bidimensionales requiere la construcción de escenarios que integren elementos visuales, distribución espacial y condiciones que permitan la navegación e interacción del jugador.

En este contexto, Rhyzome propone una solución basada en generación procedural de contenido (\textit{Procedural Content Generation}, PCG) \parencite{shaker2016}, complementada con técnicas de aprendizaje automático para apoyar la evaluación y selección de mapas.

El proyecto busca facilitar la creación de escenarios bidimensionales mediante un sistema configurable que permita generar diferentes tipos de terrenos y analizar sus características estructurales.

% ============================================================
% 2.1. DESCRIPCIÓN DEL PROBLEMA O NECESIDAD
% ============================================================
\section{Descripción del problema o necesidad}

La creación de mapas para videojuegos bidimensionales puede requerir un esfuerzo considerable de diseño, especialmente cuando es necesario producir múltiples escenarios con diferentes configuraciones, niveles de complejidad y características de navegación.

El diseño manual implica definir la distribución de los elementos del escenario, establecer zonas transitables y obstáculos, mantener la coherencia espacial y comprobar que las características del mapa correspondan con los objetivos de la experiencia de juego.

Estas actividades pueden dificultar la generación de múltiples alternativas durante las etapas de experimentación y prototipado.

La generación procedural constituye una alternativa para automatizar parcialmente este proceso, mediante algoritmos capaces de producir escenarios a partir de parámetros y reglas previamente definidos.

Sin embargo, la generación automática no garantiza por sí misma que los escenarios obtenidos presenten características estructurales adecuadas.

Un mapa generado puede contener regiones aisladas, distribuciones poco equilibradas de obstáculos o niveles de complejidad diferentes de los esperados.

Por esta razón, resulta necesario complementar la generación procedural con mecanismos que permitan analizar y evaluar los mapas obtenidos.

En respuesta a esta necesidad, Rhyzome integra algoritmos de generación procedural y modelos de aprendizaje automático, con el propósito
de apoyar la selección de escenarios a partir de características cuantificables.

De esta manera, el problema abordado no se limita a generar mapas automáticamente, sino que también considera la necesidad de evaluar sus propiedades y facilitar la selección de resultados acordes con los parámetros establecidos por el usuario.

% ============================================================
% 2.2. CONTEXTO DE APLICACION
% ============================================================

\section{Contexto de aplicación}

Rhyzome se desarrolla en el ámbito de la creación de videojuegos bidimensionales, específicamente en los procesos de generación y experimentación con escenarios virtuales.

La aplicación está orientada principalmente a desarrolladores independientes, estudiantes y personas interesadas en la creación de mapas para videojuegos.

Estos usuarios pueden requerir herramientas que faciliten la construcción de escenarios sin necesidad de diseñar manualmente cada una de sus estructuras.

El sistema contempla la generación de cuatro tipos de escenarios:

\begin{itemize}
    \item \textbf{Bosque:}
    escenarios que representan entornos naturales mediante distribuciones espaciales generadas proceduralmente.

    \item \textbf{Desierto:}
    escenarios orientados a representar terrenos de características áridas.

    \item \textbf{Ciudades:}
    escenarios que incorporan distribuciones espaciales asociadas con entornos urbanos.

    \item \textbf{Mazmorras:}
    escenarios compuestos por espacios transitables, obstáculos y regiones conectadas.
\end{itemize}

Para la generación de estos escenarios, Rhyzome utiliza dos algoritmos principales: Perlin Noise \parencite{perlin1985} y Autómatas Celulares \parencite{johnson2010}.

Perlin Noise se emplea para la generación de bosques, desiertos y ciudades, mientras que los Autómatas Celulares se utilizan para la construcción de mazmorras.

La aplicación dispone de una interfaz web que permite configurar parámetros de generación y visualizar los mapas obtenidos.

Adicionalmente, incorpora un componente de aprendizaje automático que participa en la evaluación y selección de mapas, considerando características estructurales y niveles de complejidad.

El contexto de aplicación del proyecto corresponde principalmente al prototipado, la experimentación y el apoyo al diseño
de escenarios bidimensionales.

% ============================================================
% 2.3. JUSTIFICACION TECNICA Y SOCIAL DEL USO DE IA
% ============================================================

\section{Justificación técnica y social del uso de inteligencia artificial}


La incorporación de inteligencia artificial en Rhyzome responde a la necesidad de complementar los algoritmos de generación procedural con mecanismos de evaluación que permitan analizar las características de los mapas producidos.

Desde una perspectiva técnica, los algoritmos procedurales permiten generar diferentes configuraciones espaciales mediante parámetros como dimensiones, semillas y condiciones específicas de generación.

No obstante, estos algoritmos no determinan necesariamente si un escenario generado corresponde con el nivel de complejidad
esperado por el usuario.

Por esta razón, Rhyzome incorpora modelos de aprendizaje automático supervisado que permiten clasificar características de los mapas generados.

El componente inteligente utiliza información obtenida mediante métricas estructurales, entre ellas la accesibilidad, conectividad, densidad de obstáculos y distribución de regiones transitables.

Estas características permiten representar propiedades de los escenarios y utilizarlas como entrada para los modelos de clasificación.

La integración de estos modelos busca apoyar la selección de mapas candidatos de acuerdo con niveles de complejidad previamente definidos.

De esta manera, la inteligencia artificial no reemplaza los algoritmos procedurales, sino que complementa su funcionamiento mediante un proceso de evaluación y selección.

Desde una perspectiva social y educativa, Rhyzome busca facilitar el acceso a herramientas de experimentación con generación procedural e inteligencia artificial.

La aplicación puede servir como recurso de apoyo para estudiantes y desarrolladores interesados en comprender la relación entre los parámetros de generación, las características estructurales de los mapas y los resultados obtenidos.

Asimismo, el proyecto permite explorar la aplicación de técnicas de aprendizaje automático en un contexto práctico de desarrollo de videojuegos.

La pertinencia de la solución radica, por tanto, en la integración de generación procedural, evaluación cuantitativa y aprendizaje automático dentro de una misma herramienta.

% ============================================================
% 2.4. OBJETIVO GENERAL
% ============================================================

\section{Objetivo general}

Desarrollar una aplicación web para la generación procedural de mapas bidimensionales orientados a videojuegos, mediante la implementación de los algoritmos Perlin Noise y Autómatas Celulares y la integración de modelos de aprendizaje automático que permitan evaluar características estructurales y apoyar la selección de escenarios según niveles de complejidad definidos.

% ============================================================
% 2.5. OBJETIVOS ESPECIFICOS
% ============================================================

\section{Objetivos específicos}

\begin{enumerate}

    \item Diseñar e implementar algoritmos de generación procedural basados en Perlin Noise y Autómatas Celulares para producir diferentes tipos de escenarios bidimensionales.
    
    \item Definir e implementar métricas de evaluación que permitan caracterizar las propiedades estructurales de los mapas generados, considerando aspectos como accesibilidad, conectividad y distribución de obstáculos.
    
    \item Desarrollar y evaluar modelos de aprendizaje automático supervisado para clasificar los mapas según niveles de complejidad, utilizando conjuntos de datos experimentales y métricas de desempeño.

    \item Integrar los algoritmos procedurales y el componente inteligente en una aplicación web que permita configurar parámetros, generar escenarios y visualizar los resultados obtenidos.

    \item Validar el funcionamiento del sistema mediante pruebas funcionales, de integración y evaluación de resultados,     verificando el cumplimiento de los requerimientos establecidos.
    
\end{enumerate}

